Quando um computador inicia uma aplicação, o sistema operacional aloca 
automaticamente a memória necessária para a sua execução em diferentes 
segmentos, cada um com propriedades e funções específicas.

\begin{figure}[htbp]
\centering
\includegraphics[scale=0.2]{memory-segments}
\label{fig:memorySegmentsDiagram}
\end{figure}

\subsection{Text}

A seção \texttt{Text}, contém o código executável do programa. As instruções de 
máquina presentes nesta seção são acessadas pelo processador para execução da 
aplicação. Importante destacar que esta seção é somente de leitura, o que 
significa que, uma vez definida, seu conteúdo não pode ser modificado durante 
a execução do programa. Isso é essencial para a integridade e segurança do 
código executável.

\subsection{Data}

A seção \texttt{Data}, é onde são armazenadas as variáveis globais e estáticas do 
programa que foram inicializadas durante o processo de compilação. É importante 
destacar que a seção de dados inicializados não é apenas de leitura; ela pode 
ser modificada durante a execução do programa.

Considere o seguinte código:

\begin{excode}
    static char *z = "Eu sou Z\n"
    char *w = "Eu sou W\n"

    void func(void)
    {
        static const char *y = "Eu sou Y\n"
    }

    int main(void)
    {
        const char *x = "Eu sou X\n"
        func();

        return 0;
    }
\end{excode}

No código acima, são definidas as variáveis \texttt{x}, \texttt{y}, \texttt{z} 
e \texttt{w}. Todas elas, exceto \texttt{x}, são direcionadas para a seção 
\texttt{Data}. Isso ocorre porque as variáveis globais, independentemente de 
usarem a diretiva \texttt{static}, e as variáveis locais que utilizam a 
diretiva \texttt{static}, são alocadas na seção \texttt{Data} caso sejam 
inicializadas corretamente.

\subsection{BSS}

A seção \texttt{BSS}, armazena variáveis globais e estáticas que não foram 
inicializadas explicitamente pelo programador. É normal ver também a seção 
\texttt{BSS} ser considerada uma só com a seção \texttt{Data} por ambas 
conterem variáveis globais e estáticas, diferenciando-se apenas pelo estado 
de inicialização das variáveis.

Repetindo o mesmo código da seção anterior de forma ilustrativa:

\newpage
\begin{excode}
    static char *z;
    char *w;

    void func(void)
    {
        static const char *y;
    }

    int main(void)
    {
        const char *x;
        func();

        return 0;
    }
\end{excode}

As variáveis \texttt{y}, \texttt{z} e \texttt{w}, sendo declaradas como globais 
ou estáticas sem uma inicialização explícita, são alocadas na seção \texttt{BSS} 
e inicializadas automaticamente com zero. Isso é uma característica padrão do 
compilador para assegurar que todas as variáveis na seção \texttt{BSS} comecem 
com um estado previsível.

Por outro lado, a variável \texttt{x} apesar de não inicializada é declarada 
localmente, portanto não pertence à seção \texttt{BSS}. Em C, variáveis locais 
não inicializadas como \texttt{x} podem conter valores indeterminados, pois o 
compilador não realiza nenhuma inicialização automática para elas. Isso implica 
que \texttt{x} pode conter qualquer valor aleatório presente na stack no 
momento de sua declaração.

\begin{postscript}
    \texttt{BSS} é uma abreviação de "Block Started By Symbol". 
\end{postscript}

\subsection{Heap}

A seção \texttt{Heap} é utilizada para armazenar dados alocados e liberados dinamicamente pelo programa durante sua execução, por meio de funções como \texttt{malloc}, \texttt{calloc}, \texttt{realloc} e \texttt{free}.

É importante destacar que, ao contrário das outras seções de memória, a seção \texttt{Heap} não é necessariamente contígua na memória. Ela pode utilizar técnicas de mapeamento e alocação não contígua para ajustar o tamanho dos blocos de memória alocados dinamicamente. Devido a essa maior complexidade e dinamicidade dos blocos de memória na \texttt{Heap}, acessar variáveis localizadas nessa seção pode ser mais lento em comparação com as seções estáticas.

Além disso, a seção \texttt{Heap}, conforme ilustrado na figura \ref{fig:memorySegmentsDiagram}, começa imediatamente após a seção \texttt{BSS} e pode crescer em direção à seção \texttt{Stack} à medida que mais memória dinâmica é alocada e utilizada.

%TODO: Elaborar uma melhor explicação com mais exemplos do que é a Stack
\subsection{Stack}

A seção de memória da Stack é parte crucial da memória de um programa em C, 
responsável pelo armazenamento de variáveis locais e pela manutenção do escopo 
de funções.

A Stack, também conhecida como pilha, é uma estrutura de dados de natureza 
LIFO\footnote{Last In, First Out}, o que significa que o último item a ser 
inserido é o primeiro a ser retirado. No contexto da execução de um programa em 
C, a \texttt{Stack} é utilizada para armazenar variáveis locais, parâmetros de 
função e endereços de retorno.

Quando uma função é chamada, um novo frame de execução é criado na \texttt{Stack} 
para armazenar suas variáveis locais e outros dados relacionados. Esse frame é 
desalocado quando a função retorna, liberando assim a memória alocada para suas 
variáveis locais.

É importante ressaltar que a \texttt{Stack} possui um tamanho limitado, 
determinado pelo sistema operacional ou pelo compilador. Portanto, a alocação 
excessiva de memória na \texttt{Stack} pode levar a um estouro de 
pilha\footnote{também conhecido como Stack Overflow}, resultando em falha na 
execução do programa. Além disso, como a memória na \footnote{Stack} é alocada 
de forma estática durante a compilação, não é possível alocar ou liberar memória 
de forma dinâmica, como acontece na \texttt{Heap}.

\subsubsection{Exercícios}

\textbf{É de extrema importância que você realize os exercícios aqui elaborados para fixar o conteúdo.}

\begin{enumerate}
    \item {Escreva um programa que declare uma variável global inicializada, uma 
        variável global não inicializada, uma variável local estática inicializada 
        dentro de uma função e uma variável local automática (não estática). 
        Dentro da função, modifique cada uma delas (se possível) e imprima seus 
        valores e tente explicar em qual seção de memória (Text, Data, BSS, Stack, Heap) 
        cada variável é armazenada}
    \item {Analise o código a seguir e identifique quais variáveis serão alocadas na seção \texttt{Data} e quais na seção \texttt{BSS}. Justifique sua resposta considerando se são globais/estáticas e se são inicializadas ou não.}
\newpage
    \begin{excode}
        int global_a = 10;
        int global_b;
        static float static_c = 3.14;
        static double static_d;
        
        void func(void) {
            static char static_e = 'A';
            static long static_f;
            int local_g;
        }
    \end{excode}
    \item {Considerando a figura dos segmentos de memória, descreva o que 
        acontece quando o programa solicita alocações dinâmicas muito grandes 
        através de \texttt{malloc}. Em que direção a seção \texttt{Heap} cresce? O que 
        pode ocorrer se a seção \texttt{Heap} e \texttt{Stack} se encontrarem?}
\end{enumerate}
